\documentclass[UTF8, 12pt, fontset=fandol, oneside]{ctexart}
\usepackage{researchnotes}

% 保留分册独立编译时的章号前缀。
\renewcommand{\thesection}{6.\arabic{section}}

\title{第六章\quad 不定积分}
\author{}
\date{}

\begin{document}

\maketitle

\section{原函数与不定积分}

\subsection{原函数与不定积分的概念}

大家已经知道, 寻求一个已知函数的导数是一个十分有意义的问题, 这是因为各种各样的物理问题和几何问题的解决, 常常归结为求某一函数的导数. 但是, 我们也必须看到另一方面, 那就是很多的实际问题的解决, 不是归结为求某一已知函数的导数, 而恰恰相反, 是要求一个未知函数, 使之以某一已知函数为其导数. 例如, 考虑质点沿直线运动, 已知质点所走过的路程 $s=s(t)$, 如果要求瞬时速度, 则我们有 $v(t)=s'(t)$. 反过来, 如果知道每个时刻的瞬时速度 $v(t)$, 要求质点所走过的路程 $s(t)$, 则是求导数的逆运算. 此时我们要找一个函数 $s(t)$, 使得 $s'(t)=v(t)$. 这个 $s(t)$ 称为 $v(t)$ 的一个原函数. 原函数的定义如下:

\noindent\textbf{定义 6.1.1}\quad 在区间 $I$ 上给定函数 $f(x)$, 若存在定义在 $I$ 上的函数 $F(x)$, 使得
\[
  F'(x)=f(x),\quad \forall x\in I
  \quad\text{或}\quad
  \mathrm dF(x)=f(x)\mathrm dx,\quad \forall x\in I,
\]
则称 $F(x)$ 是 $f(x)$ 的一个原函数。

由定义不难看出, 若 $f(x)$ 有原函数, 则它就有无穷多个原函数. 这是因为若 $F(x)$ 是 $f(x)$ 的一个原函数, 则对任意的常数 $C$, 函数 $F(x)+C$ 也是 $f(x)$ 的一个原函数. 几何上看这是明显的, 曲线 $F(x)$ 和 $F(x)+C$ 在点 $x$ 有相同切线斜率 (参见图 6.1.1). 另一方面, 由拉格朗日微分中值定理的推论 2 知, 若 $F(x)$ 和 $G(x)$ 是 $f(x)$ 的两个原函数, 则它们之间只能是相差一个常数, 即 $G(x)=F(x)+C$. 原函数的这一特点, 也可以从实际问题中反映出来. 例如, 两个粒子如果以相同的速度 $v(t)$ 运动, 则它们在任何相同时间走过的路程是相同的. 因此, 它们的路程函数只差一个常数, 这是由它们的初始位置的不同而引起的。

这样, 设 $f(x)$ 在区间 $I$ 上有定义, 函数 $F(x)$ 是 $f(x)$ 的一个原函数, 则函数族 $F(x)+C$ ($C$ 为任意常数) 包含了 $f(x)$ 的所有原函数。

\noindent\textbf{定义 6.1.2}\quad 若函数 $f(x)$ 在区间 $I$ 上存在原函数, 则称 $f(x)$ 的全体原函数为 $f(x)$ 的不定积分, 记为
\[
  \int f(x)\mathrm dx,
\]
这里符号 ``$\int$'' 称为积分号, $f(x)$ 称为被积函数, $x$ 称为积分变量。

根据上述定义以及上面的分析我们有, 若 $F(x)$ 是 $f(x)$ 的一个原函数, 则
\[
  \int f(x)\mathrm dx=F(x)+C,
\]
其中 $C$ 为任意常数. 另外我们有
\[
  \left(\int f(x)\mathrm dx\right)'=f(x),
  \qquad
  \mathrm d\left(\int f(x)\mathrm dx\right)=f(x)\mathrm dx
\]
和
\[
  \int F'(x)\mathrm dx=F(x)+C,
  \qquad
  \int \mathrm dF(x)=F(x)+C,
\]
其中 $C$ 为任意常数. 因此, 在允许相差一个常数的意义下, 求不定积分的这一运算恰好是求导或求微分的逆运算。

\subsection{基本不定积分表和不定积分的线性性质}

由于微分和积分互为逆运算, 因此对应于微分的每一种方法相应地就有一种积分的方法. 首先, 从求导数的公式中我们可以导出如下不定积分公式, 它是我们计算不定积分的基础, 务必牢记。
\[
\begin{aligned}
  &(1)\quad \int x^\alpha\mathrm dx=\frac1{\alpha+1}x^{\alpha+1}+C
  \quad(\alpha\ne-1);\\
  &(2)\quad \int\frac{\mathrm dx}{x}=\ln|x|+C;\\
  &(3)\quad \int e^x\mathrm dx=e^x+C;\\
  &(4)\quad \int\cos x\mathrm dx=\sin x+C;\\
  &(5)\quad \int\sin x\mathrm dx=-\cos x+C;\\
  &(6)\quad \int\frac{\mathrm dx}{\cos^2x}=\tan x+C;\\
  &(7)\quad \int\frac{\mathrm dx}{1+x^2}=\arctan x+C;\\
  &(8)\quad \int\frac{\mathrm dx}{\sqrt{1-x^2}}=\arcsin x+C;\\
  &(9)\quad \int\sinh x\mathrm dx=\cosh x+C;\\
  &(10)\quad \int\cosh x\mathrm dx=\sinh x+C;\\
  &(11)\quad \int\frac{\mathrm dx}{\cosh^2x}=\tanh x+C.
\end{aligned}
\]
其次, 对应于求导法则
\[
  (f(x)\pm g(x))'=f'(x)\pm g'(x),
\]
就有如下的求不定积分性质:

\noindent\textbf{性质 6.1.1}\quad 设函数 $f(x),g(x)$ 存在原函数, 则 $f(x)\pm g(x)$ 也存在原函数, 而且有
\[
  \int[f(x)\pm g(x)]\mathrm dx
  =
  \int f(x)\mathrm dx\pm\int g(x)\mathrm dx.
\]

\noindent\textbf{证明}\quad 令
\[
  \int f(x)\mathrm dx=F(x)+C
  \quad\text{或}\quad
  F'(x)=f(x),
\]
\[
  \int g(x)\mathrm dx=G(x)+C
  \quad\text{或}\quad
  G'(x)=g(x),
\]
则
\[
  [F(x)\pm G(x)]'=f(x)\pm g(x),
\]
所以
\[
  \int[f(x)\pm g(x)]\mathrm dx=F(x)\pm G(x)+C.
\]

对应于求导法则 $(kf(x))'=kf'(x)$, 其中 $k$ 是常数, 就有如下的求不定积分性质:

\noindent\textbf{性质 6.1.2}\quad 设函数 $f(x)$ 存在原函数, 则 $kf(x)$ 存在原函数, 且
\[
  \int kf(x)\mathrm dx=k\int f(x)\mathrm dx,
\]
其中 $k$ 是常数, 且 $k\ne0$。

\noindent\textbf{证明留作练习.}

性质 6.1.1 和 6.1.2 说明求不定积分运算是一种线性运算。

\noindent\textbf{例 6.1.1}\quad 求不定积分
\[
  \int\frac{x^4}{1+x^2}\mathrm dx.
\]

\noindent\textbf{解}\quad 将被积函数恒等变形, 有
\[
\begin{aligned}
  \int\frac{x^4}{1+x^2}\mathrm dx
  &=
  \int\frac{x^4-1+1}{1+x^2}\mathrm dx\\
  &=
  \int\left(x^2-1+\frac1{1+x^2}\right)\mathrm dx\\
  &=
  \frac{x^3}{3}-x+\arctan x+C.
\end{aligned}
\]

\noindent\textbf{例 6.1.2}\quad 求不定积分
\[
  \int\tan^2x\mathrm dx.
\]

\noindent\textbf{解}\quad 将被积函数恒等变形, 有
\[
\begin{aligned}
  \int\tan^2x\mathrm dx
  &=
  \int\frac{\sin^2x}{\cos^2x}\mathrm dx\\
  &=
  \int\left(\frac1{\cos^2x}-1\right)\mathrm dx
  =
  \tan x-x+C.
\end{aligned}
\]

\noindent\textbf{例 6.1.3}\quad 求不定积分
\[
  \int\frac1{\sin^22x}\mathrm dx.
\]

\noindent\textbf{解}\quad
\[
\begin{aligned}
  \int\frac1{\sin^22x}\mathrm dx
  &=
  \int\frac1{4\sin^2x\cos^2x}\mathrm dx\\
  &=
  \int\frac{\sin^2x+\cos^2x}{4\sin^2x\cos^2x}\mathrm dx\\
  &=
  \frac14\left(\int\frac1{\cos^2x}\mathrm dx+\int\frac1{\sin^2x}\mathrm dx\right)\\
  &=
  \frac14(\tan x-\cot x)+C.
\end{aligned}
\]

\section{换元法与分部积分法}

不定积分换元法是对应于复合函数的求导法则的. 大家已经知道, 对于复合函数而言, 我们有如下的一阶微分形式不变性
\[
  f(u(x))u'(x)\mathrm dx=f(u)\mathrm du,
\]
其中 $u=u(x)$. 由此立得
\begin{equation}
  \int f(u(x))u'(x)\mathrm dx=\int f(u)\mathrm du.
  \tag{6.2.1}
\end{equation}

利用这一等式可导出两种求不定积分的方法: 如果所求的不定积分有 (6.2.1) 的左边所示的形式, 则我们可以利用变量替换 $u=u(x)$ 将其转化为 (6.2.1) 右边的形式, 然后通过求右边关于变量 $u$ 的不定积分而得到所要求的不定积分, 这就是所谓的\textbf{第一换元法}; 如果所求的不定积分有 (6.2.1) 的右边所示的形式, 则我们可以利用变量替换 $u=u(x)$ 将其转化为 (6.2.1) 左边的形式, 然后通过求左边关于变量 $x$ 的不定积分而得到所要求的不定积分, 这就是所谓的\textbf{第二换元法}。

\subsection{第一换元法}

\noindent\textbf{定理 6.2.1}\quad 如果
\[
  \int f(u)\mathrm du=F(u)+C,
\]
而 $u=u(x)$ 是关于 $x$ 的可微函数, 则有
\[
  \int f(u(x))u'(x)\mathrm dx=F(u(x))+C.
\]

\noindent\textbf{证明}\quad 由定理的条件, 我们有
\[
  \mathrm dF(u)=f(u)\mathrm du.
\]
再根据一阶微分有形式不变性, 可得
\[
  \mathrm dF(u(x))=f(u(x))\mathrm du(x)=f(u(x))u'(x)\mathrm dx,
\]
所以
\[
  \int f(u(x))u'(x)\mathrm dx=F(u(x))+C.
\]

\noindent\textbf{例 6.2.1}\quad 求不定积分
\[
  \int\frac1{x-1}\mathrm dx.
\]

\noindent\textbf{解}\quad 我们已经知道
\[
  \int\frac{\mathrm du}{u}=\ln|u|+C.
\]
如果令 $u=x-1$, 则 $u$ 是关于 $x$ 的可微函数, 而且 $\mathrm du=\mathrm dx$, 因此我们有
\[
  \int\frac1{x-1}\mathrm dx
  =
  \int\frac{\mathrm du}{u}
  =
  \ln|u|+C
  =
  \ln|x-1|+C.
\]

\noindent\textbf{例 6.2.2}\quad 求不定积分
\[
  \int\frac{\mathrm dx}{ax+b}\quad (a\ne0).
\]

\noindent\textbf{解}\quad 令 $u=ax+b$, 则有 $\mathrm du=a\mathrm dx$, 即 $\mathrm dx=\dfrac{\mathrm du}{a}$, 于是有
\[
  \int\frac{\mathrm dx}{ax+b}
  =
  \int\frac1a\frac{\mathrm du}{u}
  =
  \frac1a\ln|u|+C
  =
  \frac1a\ln|ax+b|+C.
\]

\noindent\textbf{例 6.2.3}\quad 求不定积分
\[
  \int(ax+b)^n\mathrm dx\quad (a\ne0,\ n\ne-1).
\]

\noindent\textbf{解}\quad 令 $u=ax+b$, 则有
\[
  \int(ax+b)^n\mathrm dx
  =
  \frac1a\int u^n\mathrm du
  =
  \frac{u^{n+1}}{a(n+1)}+C
  =
  \frac{(ax+b)^{n+1}}{a(n+1)}+C.
\]

\noindent\textbf{例 6.2.4}\quad 求不定积分
\[
  \int\frac{\mathrm dx}{a^2+x^2}\quad (a\ne0).
\]

\noindent\textbf{解}\quad 由于
\[
  \int\frac{\mathrm dx}{a^2+x^2}
  =
  \frac1{a^2}\int\frac{\mathrm dx}{1+\left(\dfrac xa\right)^2},
\]
因此令 $u=\dfrac xa$, 则有 $\mathrm du=\dfrac{\mathrm dx}{a}$, 即 $\mathrm dx=a\mathrm du$, 从而
\[
  \int\frac{\mathrm dx}{a^2+x^2}
  =
  \frac1a\int\frac{\mathrm du}{1+u^2}
  =
  \frac1a\arctan\frac xa+C.
\]

\noindent\textbf{例 6.2.5}\quad 求不定积分
\[
  \int\frac{\mathrm dx}{x^4-1}.
\]

\noindent\textbf{解}\quad 将被积函数恒等变形, 有
\[
\begin{aligned}
  \int\frac{\mathrm dx}{x^4-1}
  &=
  \int\frac{\mathrm dx}{(x^2-1)(x^2+1)}\\
  &=
  \frac12\left(\int\frac{\mathrm dx}{x^2-1}-\int\frac{\mathrm dx}{x^2+1}\right)\\
  &=
  \frac12\left[\frac12\left(\int\frac{\mathrm dx}{x-1}-\int\frac{\mathrm dx}{x+1}\right)
  -\int\frac{\mathrm dx}{x^2+1}\right]\\
  &=
  \frac14\ln\left|\frac{x-1}{x+1}\right|-\frac12\arctan x+C.
\end{aligned}
\]

从以上例子可以看出, 采用第一换元法求不定积分的关键是, 要将所求不定积分的被积表达式 $f(x)\mathrm dx$ 写成两个因式的乘积, 前一因式形如 $g(u(x))$, 后一因式形如 $\mathrm du(x)$. 对于与上面例子类似的一些简单的不定积分, 我们很容易做到这一点. 但对于较为复杂的不定积分, 要做到这一点就比较困难. 要对具体问题做具体分析, 灵活运用. 下面看一些较为复杂的例子。

\noindent\textbf{例 6.2.6}\quad 求不定积分
\[
  \int\frac{x\mathrm dx}{x^2+1}.
\]

\noindent\textbf{解}\quad 令 $u=x^2+1$, 则有 $\mathrm du=2x\mathrm dx$, 即 $x\mathrm dx=\dfrac12\mathrm du$, 所以
\[
  \int\frac{x\mathrm dx}{x^2+1}
  =
  \frac12\int\frac{\mathrm d(x^2+1)}{x^2+1}
  =
  \frac12\ln(x^2+1)+C.
\]
在此例中, 我们在计算过程中没有写出关于 $u$ 的积分. 在熟练后, 我们不必每次写出代表中间变量的符号 $u$, 而只要在心目中将变量 $u$ 所代表的函数 $u(x)$ 看做一个整体就可以了. 此外, 从这个例子也可以看出, 确定 $u=u(x)$ 的时候, 要仔细观察在除去因子 $g(u(x))$ 之后, 是否剩下的函数恰好为 $u'(x)$ (或相差一个常数倍). 这些技巧读者应该在不断的练习过程中加以总结和掌握。

\noindent\textbf{例 6.2.7}\quad 求不定积分
\[
  \int\frac{\mathrm dx}{a^2-x^2}\quad (a\ne0).
\]

\noindent\textbf{解法 1}\quad
\[
\begin{aligned}
  \int\frac{\mathrm dx}{a^2-x^2}
  &=
  \frac1a\int\frac{\mathrm d\left(\dfrac xa\right)}{1-\left(\dfrac xa\right)^2}\\
  &=
  \frac1{2a}\ln\left|\frac{1+\dfrac xa}{1-\dfrac xa}\right|+C
  =
  \frac1{2a}\ln\left|\frac{a+x}{a-x}\right|+C.
\end{aligned}
\]

\noindent\textbf{解法 2}\quad
\[
\begin{aligned}
  \int\frac{\mathrm dx}{a^2-x^2}
  &=
  \frac1{2a}\int\left(\frac1{a+x}+\frac1{a-x}\right)\mathrm dx\\
  &=
  \frac1{2a}\left[\int\frac{\mathrm d(a+x)}{a+x}
  -\int\frac{\mathrm d(a-x)}{a-x}\right]\\
  &=
  \frac1{2a}\ln\left|\frac{a+x}{a-x}\right|+C.
\end{aligned}
\]

\noindent\textbf{例 6.2.8}\quad 求不定积分
\[
  I=\int\frac{\mathrm dx}{\sqrt{a^2-x^2}}\quad (a>0).
\]

\noindent\textbf{解}\quad 将被积函数变形得
\[
  I=
  \int\frac{\mathrm d\left(\dfrac xa\right)}
  {\sqrt{1-\left(\dfrac xa\right)^2}}
  =
  \arcsin\frac xa+C.
\]

\noindent\textbf{例 6.2.9}\quad 求不定积分
\[
  I=\int\frac{\mathrm dx}{\sqrt{x(1-x)}}.
\]

\noindent\textbf{解法 1}\quad 将被积函数变形得
\[
  I=
  \int\frac{\mathrm dx}{\sqrt{\dfrac14-\left(x-\dfrac12\right)^2}}
  =
  \arcsin(2x-1)+C.
\]

\noindent\textbf{解法 2}\quad 由 $\mathrm dx/\sqrt x=2\mathrm d\sqrt x$, 有
\[
  I=
  \int\frac{\mathrm dx}{\sqrt x\sqrt{1-x}}
  =
  2\int\frac{\mathrm d\sqrt x}{\sqrt{1-(\sqrt x)^2}}
  =
  2\arcsin(\sqrt x)+C.
\]

\noindent\textbf{注}\quad 虽然两种方法得到的答案形式不同, 但实质上它们只相差一个常数. 这一点请读者自行验证之。

\noindent\textbf{例 6.2.10}\quad 求不定积分
\[
  I=\int\tan x\mathrm dx.
\]

\noindent\textbf{解}\quad 利用 $\sin x\mathrm dx=-\mathrm d\cos x$, 有
\[
  I=-\int\frac{\mathrm d\cos x}{\cos x}
  =
  -\ln|\cos x|+C.
\]

\noindent\textbf{例 6.2.11}\quad 求不定积分
\[
  I_n=\int\tan^nx\mathrm dx.
\]

\noindent\textbf{解}\quad 利用 $\tan^2x=\sec^2x-1$, 有
\[
\begin{aligned}
  I_n
  &=
  \int\tan^{n-2}x\frac{1-\cos^2x}{\cos^2x}\mathrm dx\\
  &=
  \int\tan^{n-2}x\mathrm d\tan x-\int\tan^{n-2}x\mathrm dx\\
  &=
  \frac1{n-1}\tan^{n-1}x-I_{n-2}.
\end{aligned}
\]
这是一个递推公式. 利用这一公式, 我们可以归纳地求出 $I_n$. 例如
\[
  I_3=\frac12\tan^2x-\int\tan x\mathrm dx
  =
  \frac12\tan^2x+\ln|\cos x|+C;
\]
\[
\begin{aligned}
  I_4
  &=
  \frac13\tan^3x-\int\tan^2x\mathrm dx\\
  &=
  \frac13\tan^3x-\tan x+\int1\mathrm dx\\
  &=
  \frac13\tan^3x-\tan x+x+C.
\end{aligned}
\]

\noindent\textbf{例 6.2.12}\quad 求不定积分
\[
  \int\frac{\mathrm dx}{\cos x}.
\]

\noindent\textbf{解法 1}\quad 将被积函数变形, 并利用 $\cos x\mathrm dx=\mathrm d\sin x$ 得
\[
\begin{aligned}
  I
  &=
  \int\frac{\cos x\mathrm dx}{\cos^2x}
  =
  \int\frac{\mathrm d(\sin x)}{1-\sin^2x}\\
  &=
  \frac12\ln\left|\frac{1+\sin x}{1-\sin x}\right|+C
  =
  \frac12\ln\left|\frac{1+\sin x}{\cos x}\right|^2+C\\
  &=
  \ln|\sec x+\tan x|+C.
\end{aligned}
\]

\noindent\textbf{解法 2}\quad 将被积函数变形得
\[
\begin{aligned}
  I
  &=
  \int\frac{\mathrm dx}{\cos^2\dfrac x2-\sin^2\dfrac x2}
  =
  2\int\frac{\mathrm d\left(\dfrac x2\right)}
  {\cos^2\dfrac x2\left(1-\tan^2\dfrac x2\right)}\\
  &=
  2\int\frac{\mathrm d\tan\left(\dfrac x2\right)}
  {1-\tan^2\dfrac x2}
  =
  \ln\left|\frac{1+\tan\dfrac x2}{1-\tan\dfrac x2}\right|+C\\
  &=
  \ln\left|\tan\left(\frac x2+\frac\pi4\right)\right|+C.
\end{aligned}
\]

\noindent\textbf{例 6.2.13}\quad 求不定积分
\[
  I=\int\frac{\mathrm dx}{1+x^3}
  \quad\text{和}\quad
  J=\int\frac{x\mathrm dx}{1+x^3}.
\]

\noindent\textbf{解}\quad 由于
\[
\begin{aligned}
  I+J
  &=
  \int\frac{1+x}{1+x^3}\mathrm dx
  =
  \int\frac{\mathrm dx}{1-x+x^2}
  =
  \int\frac{\mathrm dx}{\left(x-\dfrac12\right)^2+\dfrac34}\\
  &=
  \frac2{\sqrt3}\arctan\frac{2x-1}{\sqrt3}+C,
\end{aligned}
\]
\[
\begin{aligned}
  I-J
  &=
  \int\frac{1-x}{1+x^3}\mathrm dx
  =
  \int\frac{\mathrm dx}{1+x}
  -
  \int\frac{x^2\mathrm dx}{1+x^3}\\
  &=
  \ln|1+x|-\frac13\ln|1+x^3|+C,
\end{aligned}
\]
所以
\[
  I=
  \frac1{\sqrt3}\arctan\frac{2x-1}{\sqrt3}
  +
  \frac12\ln|1+x|-\frac16\ln|1+x^3|+C;
\]
\[
  J=
  \frac1{\sqrt3}\arctan\frac{2x-1}{\sqrt3}
  -
  \frac12\ln|1+x|+\frac16\ln|1+x^3|+C.
\]

\subsection{第二换元法}

\noindent\textbf{定理 6.2.2}\quad 设函数 $x=x(t)$ 在某一开区间上可导, 且 $x'(t)\ne0$. 如果
\[
  \int f(x(t))x'(t)\mathrm dt=G(t)+C,
\]
则有
\[
  \int f(x)\mathrm dx=G(t(x))+C,
\]
其中 $t=t(x)$ 为 $x=x(t)$ 的反函数。

\noindent\textbf{证明}\quad 由不定积分的定义, 可得 $G'(t)=f(x(t))x'(t)$. 再由 $x'(t)\ne0$ 和达布定理知, $x'(t)$ 恒大于 $0$ 或恒小于 $0$. 所以 $x(t)$ 是严格单调连续函数, 从而其反函数 $t=t(x)$ 存在、连续、且严格单调. 注意到
\[
  t'(x)=\frac1{x'(t(x))},
\]
便有
\[
  (G(t(x)))'=G'(t(x))t'(x)=f(x)x'(t(x))t'(x)=f(x),
\]
所以
\[
  \int f(x)\mathrm dx=G(t(x))+C.
\]

第二换元法主要用来求含有根式的不定积分。

\noindent\textbf{例 6.2.14}\quad 求不定积分
\[
  I=\int\sqrt{a^2-x^2}\mathrm dx.
\]

\noindent\textbf{解}\quad 令 $x=a\sin t,\ |t|<\pi/2$, 则
\[
  I=\int a^2\cos^2t\mathrm dt
  =
  \frac{a^2}2t+\frac{a^2}4\sin2t+C
  =
  \frac{a^2}2\arcsin\frac xa+\frac12x\sqrt{a^2-x^2}+C.
\]
在将上面积分结果中的新变量 $t$ 返回到变量 $x$ 时, 我们根据变换 $x=a\sin t$ 作直角三角形 (见图 6.2.1), 可得出
\[
  \cos t=\frac{\sqrt{a^2-x^2}}{a}.
\]

\noindent\textbf{例 6.2.15}\quad 求不定积分
\[
  I=\int\frac1{\sqrt{x^2-a^2}}\mathrm dx\quad (a>0,\ |x|>a).
\]

\noindent\textbf{解法 1}\quad 令 $x=a\sec t\ (0<t<\pi/2)$, 则
\[
\begin{aligned}
  I
  &=
  \int\frac{a\sec t\tan t}{a\tan t}\mathrm dt
  =
  \int\frac{\mathrm dt}{\cos t}
  =
  \ln|\sec t+\tan t|+C\\
  &=
  \ln\left|\frac xa+\frac{\sqrt{x^2-a^2}}a\right|+C
  =
  \ln|x+\sqrt{x^2-a^2}|+C.
\end{aligned}
\]
在将上面积分结果中的新变量 $t$ 返回到变量 $x$ 时, 我们利用了图 6.2.2, 其中 $\sec t=x/a,\ \tan t=\sqrt{x^2-a^2}/a$。

\noindent\textbf{解法 2}\quad 令 $x=a\cosh t,\ 0<t<+\infty$, 则 $t=\ln|x+\sqrt{x^2-a^2}|-\ln a$, 而且
\[
  I=\int\frac{a\sinh t}{a\sinh t}\mathrm dt
  =
  t+C
  =
  \ln|x+\sqrt{x^2-a^2}|+C.
\]

\noindent\textbf{注}\quad 上面是对 $x>a$ 进行积分的, 对于 $x<-a$ 情形, 可以用同样方法处理。

\noindent\textbf{例 6.2.16}\quad 求不定积分
\[
  I=\int\frac1{\sqrt{x^2+a^2}}\mathrm dx\quad (a>0).
\]

\noindent\textbf{解法 1}\quad 令 $x=a\sinh t$, 类似于上面的解法 2, 我们有
\[
  I=t+C=\operatorname{arcsinh}\frac xa+C
  =
  \ln(x+\sqrt{x^2+a^2})+C.
\]

\noindent\textbf{解法 2}\quad 令 $x=a\tan t$, 则 $\mathrm dx=a\sec^2t\mathrm dt$,
\[
\begin{aligned}
  I
  &=
  \int\frac1{a\sec t}\cdot a\sec^2t\mathrm dt
  =
  \int\frac1{\cos t}\mathrm dt\\
  &=
  \ln|\sec t+\tan t|+C
  =
  \ln\left|\frac{\sqrt{a^2+x^2}}a+\frac xa\right|+C\\
  &=
  \ln|x+\sqrt{a^2+x^2}|+C.
\end{aligned}
\]
将 $t$ 还原 $x$ 的过程中, 用到图 6.2.3 中 $\tan t=x/a$ 的关系。

\noindent\textbf{例 6.2.17}\quad 求不定积分
\[
  I=\int\frac1{\sqrt x+\sqrt[3]x}\mathrm dx.
\]

\noindent\textbf{解}\quad 令 $x=t^6$, 则有
\[
\begin{aligned}
  I
  &=
  \int\frac{6t^5\mathrm dt}{t^3+t^2}
  =
  6\int\frac{t^3}{1+t}\mathrm dt\\
  &=
  6\int\left(t^2-t+1-\frac1{1+t}\right)\mathrm dt\\
  &=
  2t^3-3t^2+6t-6\ln|1+t|+C\\
  &=
  2\sqrt x-3\sqrt[3]x+6\sqrt[6]x-6\ln|1+\sqrt[6]x|+C.
\end{aligned}
\]

\subsection{分部积分法}

对应于求导法则
\[
  (u(x)v(x))'=u'(x)v(x)+u(x)v'(x),
\]
有如下的分部积分法:

\noindent\textbf{定理 6.2.3}\quad 设 $u(x),v(x)$ 可导, 若 $\int u'(x)v(x)\mathrm dx$ 存在, 则
\[
  \int u(x)v'(x)\mathrm dx
  =
  u(x)v(x)-\int u'(x)v(x)\mathrm dx.
\]

\noindent\textbf{证明}\quad 由 $(uv)'=u'v+uv'$, 我们有 $uv'=(uv)'-u'v$. 而该等式的右端两项的原函数都存在, 故其左端项的原函数也存在, 且有
\[
  \int uv'\mathrm dx=uv-\int u'v\mathrm dx.
\]

\noindent\textbf{注}\quad 为了便于记忆, 常将分部积分公式写成
\[
  \int u\mathrm dv=uv-\int v\mathrm du.
\]
用分部积分法求不定积分的一般步骤为:

(1) 把被积函数拆成 $uv'$, 并且将 $v'\mathrm dx$ 改写成 $\mathrm dv$, 通常 $v'$ 取
\[
  e^{\pm x},\ \sin x,\ \cos x,\ x^n,\ \sinh x,\ \cosh x
\]
等函数;

(2) 使用分部积分公式;

(3) 求不定积分 $\int u'v\mathrm dx$, 或者设法建立函数方程来求解。

\noindent\textbf{例 6.2.18}\quad 求不定积分
\[
  I=\int x^3\ln x\mathrm dx.
\]

\noindent\textbf{解}\quad 令 $u=\ln x,\ \mathrm dv=x^3\mathrm dx=\mathrm d\dfrac{x^4}4$, 即 $v=\dfrac{x^4}4$, 则
\[
  I=\frac14x^4\ln x-\frac14\int x^3\mathrm dx
  =
  \frac14x^4\ln x-\frac1{16}x^4+C.
\]

\noindent\textbf{例 6.2.19}\quad 求不定积分
\[
  I=\int\arctan x\mathrm dx.
\]

\noindent\textbf{解}\quad 设 $v=x,\ u=\arctan x$, 则
\[
  I
  =
  x\arctan x-\int x\mathrm d(\arctan x)
  =
  x\arctan x-\int\frac{x\mathrm dx}{1+x^2}
  =
  x\arctan x-\frac12\ln(1+x^2)+C.
\]

\noindent\textbf{例 6.2.20}\quad 求不定积分
\[
  I=\int x^2\sin x\mathrm dx.
\]

\noindent\textbf{解}\quad 设 $v=\cos x,\ u=x^2$, 即 $\sin x\mathrm dx=-\mathrm dv$, 则
\[
\begin{aligned}
  I
  &=
  -\int x^2\mathrm d(\cos x)
  =
  -x^2\cos x+\int\cos x\mathrm d x^2\\
  &=
  -x^2\cos x+2\int x\mathrm d(\sin x)\\
  &=
  -x^2\cos x+2x\sin x-2\int\sin x\mathrm dx\\
  &=
  -x^2\cos x+2x\sin x+2\cos x+C.
\end{aligned}
\]

\noindent\textbf{例 6.2.21}\quad 求不定积分
\[
  I=\int\sqrt{x^2-1}\mathrm dx.
\]

\noindent\textbf{解法 1}\quad 由于
\[
\begin{aligned}
  I
  &=
  x\sqrt{x^2-1}-\int x\mathrm d(\sqrt{x^2-1})\\
  &=
  x\sqrt{x^2-1}-\int\frac{x^2}{\sqrt{x^2-1}}\mathrm dx\\
  &=
  x\sqrt{x^2-1}-\int\sqrt{x^2-1}\mathrm dx-\int\frac{\mathrm dx}{\sqrt{x^2-1}}\\
  &=
  x\sqrt{x^2-1}-\ln|x+\sqrt{x^2-1}|-I,
\end{aligned}
\]
所以
\[
  I=\frac12x\sqrt{x^2-1}-\frac12\ln|x+\sqrt{x^2-1}|+C.
\]

\noindent\textbf{解法 2}\quad 令 $x=\cosh t$, 则有
\[
\begin{aligned}
  I
  &=
  \int\sinh^2t\mathrm dt
  =
  \int\frac{\cosh2t-1}{2}\mathrm dt\\
  &=
  \frac14\sinh2t-\frac12t+C\\
  &=
  \frac12x\sqrt{x^2-1}-\frac12\ln|x+\sqrt{x^2-1}|+C.
\end{aligned}
\]

\noindent\textbf{例 6.2.22}\quad 求不定积分
\[
  I=\int e^{ax}\cos bx\mathrm dx
  \quad\text{和}\quad
  J=\int e^{ax}\sin bx\mathrm dx.
\]

\noindent\textbf{解}\quad 利用分部积分公式, 可得
\[
  I=\frac1b\int e^{ax}\mathrm d(\sin bx)
  =
  \frac1be^{ax}\sin bx-\frac abJ;
\]
\[
  J=-\frac1b\int e^{ax}\mathrm d(\cos bx)
  =
  -\frac1be^{ax}\cos bx+\frac abI.
\]
所以
\[
  \begin{cases}
    bI+aJ=e^{ax}\sin bx,\\
    aI-bJ=e^{ax}\cos bx.
  \end{cases}
\]
解此方程组, 可得
\[
  I=\frac{e^{ax}(a\cos bx+b\sin bx)}{a^2+b^2}+C,
\]
\[
  J=\frac{e^{ax}(a\sin bx-b\cos bx)}{a^2+b^2}+C.
\]

\noindent\textbf{例 6.2.23}\quad 求不定积分
\[
  I=\int\frac{\sin x}{a\cos x+b\sin x}\mathrm dx
  \quad\text{和}\quad
  J=\int\frac{\cos x}{a\cos x+b\sin x}\mathrm dx.
\]

\noindent\textbf{解}\quad 由于
\[
  bI+aJ=
  \int\frac{a\cos x+b\sin x}{a\cos x+b\sin x}\mathrm dx
  =
  x+C,
\]
\[
  bJ-aI=
  \int\frac{\mathrm d(a\cos x+b\sin x)}{a\cos x+b\sin x}
  =
  \ln|a\cos x+b\sin x|+C,
\]
所以
\[
  I=\frac1{a^2+b^2}(bx-a\ln|a\cos x+b\sin x|)+C,
\]
\[
  J=\frac1{a^2+b^2}(ax+b\ln|a\cos x+b\sin x|)+C.
\]

\noindent\textbf{例 6.2.24}\quad 求不定积分
\[
  K_n=\int\cos^nx\mathrm dx.
\]

\noindent\textbf{解}\quad 由于
\[
\begin{aligned}
  K_n
  &=
  \int\cos^{n-1}x\mathrm d(\sin x)
  =
  \sin x\cos^{n-1}x-\int\sin x\mathrm d(\cos^{n-1}x)\\
  &=
  \sin x\cos^{n-1}x+(n-1)\int\sin^2x\cos^{n-2}x\mathrm dx\\
  &=
  \sin x\cos^{n-1}x+(n-1)\int\cos^{n-2}x\mathrm dx
  -(n-1)\int\cos^nx\mathrm dx\\
  &=
  \sin x\cos^{n-1}x+(n-1)K_{n-2}-(n-1)K_n,
\end{aligned}
\]
所以
\[
  K_n=\frac1n\sin x\cos^{n-1}x+\frac{n-1}nK_{n-2}.
\]
由此可推出:
\[
  K_0=x+C;\quad
  K_1=\sin x+C;\quad
  K_2=\frac14\sin2x+\frac12x+C;\quad
  \cdots\cdots\cdots
\]

\noindent\textbf{例 6.2.25}\quad 求不定积分
\[
  I_n=\int\sin^nx\mathrm dx.
\]

\noindent\textbf{解}\quad 由于
\[
\begin{aligned}
  I_n
  &=
  -\int\sin^{n-1}x\mathrm d\cos x\\
  &=
  -\sin^{n-1}x\cos x+(n-1)\int\cos^2x\sin^{n-2}x\mathrm dx\\
  &=
  -\sin^{n-1}x\cos x+(n-1)I_{n-2}-(n-1)I_n,
\end{aligned}
\]
所以
\[
  I_n=-\frac{\sin^{n-1}x\cos x}{n}+\frac{n-1}{n}I_{n-2}.
\]
再注意到
\[
  I_0=x+C,\qquad I_1=-\cos x+C,
\]
便有
\[
  I_2=-\frac12\sin x\cos x+\frac12x+C;
\]
\[
  I_3=-\frac13\sin^2x\cos x-\frac23\cos x+C;
\]
\[
  I_4=-\frac14\sin^3x\cos x-\frac38\sin x\cos x+\frac38x+C;
\]
\[
  I_5=-\frac15\sin^4x\cos x-\frac4{15}\sin^2x\cos x-\frac8{15}\cos x+C;
\]
\[
  \cdots\cdots\cdots
\]

\noindent\textbf{例 6.2.26}\quad 求不定积分
\[
  I_n=\int\frac{\mathrm dx}{(x^2+a^2)^n}.
\]

\noindent\textbf{解}\quad 由于
\[
\begin{aligned}
  I_n
  &=
  \frac{x}{(x^2+a^2)^n}
  +
  2n\int\frac{x^2}{(x^2+a^2)^{n+1}}\mathrm dx\\
  &=
  \frac{x}{(x^2+a^2)^n}
  +
  2nI_n-2na^2I_{n+1},
\end{aligned}
\]
所以
\[
  I_{n+1}=
  \frac{x}{2na^2(x^2+a^2)^n}+\frac{2n-1}{2na^2}I_n.
\]
特别地, 我们有
\[
  I_2=
  \int\frac{\mathrm dx}{(x^2+a^2)^2}
  =
  \frac{x}{2a^2(x^2+a^2)}+\frac1{2a^3}\arctan\frac xa+C.
\]

从理论上来讲, 任何初等函数都有原函数 (见下章的定积分理论). 上面我们所举的例子中遇到的初等函数的原函数均为初等函数, 但并非所有的初等函数的原函数都是初等函数. 例如, 已经证明如下的积分就都不能表成初等函数:
\[
  \int e^{-x^2}\mathrm dx,\qquad
  \int\frac{\sin x}{x}\mathrm dx,\qquad
  \int\frac{\cos x}{x}\mathrm dx,
\]
\[
  \int\frac1{\ln x}\mathrm dx,\qquad
  \int\sin x^2\mathrm dx,\qquad
  \int\cos x^2\mathrm dx.
\]
再如, 积分 $\int(a+bx)^px^q\mathrm dx$ 和椭圆积分
\[
  \int R(x,\sqrt{ax^3+bx^2+cx+d})\mathrm dx
\]
以及
\[
  \int R(x,\sqrt{ax^4+bx^3+cx^2+dx+e})\mathrm dx
\]
也都不能用初等函数表示, 其中 $p,q$ 和 $p+q$ 非整数, $R$ 为两个变元的有理函数。

\section{其他类型函数的不定积分}

\subsection{有理函数的不定积分}

\textbf{有理函数}是指两个多项式之比, 即有理函数有如下形式
\[
  R(x)=\frac{P(x)}{Q(x)},
\]
其中 $P(x),Q(x)$ 为互素多项式. 理论上讲, 它的不定积分一定能够表成初等函数。

有理函数可分为真分式和假分式: \textbf{真分式}是指分子次数小于分母次数的有理函数; \textbf{假分式}是指分子次数大于或等于分母次数的有理函数. 显然, 通过作多项式除法, 任何一个假分式都可以写成一个多项式与一个真分式之和. 因此我们只需讨论真分式的不定积分即可。

从理论上来讲, 任何一个真分式都可以写成最简真分式之和. \textbf{最简真分式}是指分母为素多项式或素多项式之幂的真分式. 对于系数均为实数的多项式来讲, 素多项式只有两种, 即 $x-a$ 和 $x^2+px+q$, 其中 $p^2-4q<0$. 因此, 最简真分式只有如下四种:
\[
  \frac A{x-a},\qquad
  \frac A{(x-a)^m}\ (m>1),
\]
\[
  \frac{Ax+B}{x^2+px+q},\qquad
  \frac{Ax+B}{(x^2+px+q)^n}\ (n>1),
\]
其中 $p^2-4q<0$. 而前面所介绍的不定积分方法和例子表明, 我们已经有办法把这些最简真分式的原函数求出来, 而且它们的原函数均为初等函数。

事实上, 我们有
\[
  (1)\quad
  \int\frac A{x-a}\mathrm dx
  =
  A\int\frac1{x-a}\mathrm d(x-a)
  =
  A\ln|x-a|+C.
\]
\[
  (2)\quad
  \int\frac A{(x-a)^m}\mathrm dx
  =
  A\int\frac1{(x-a)^m}\mathrm d(x-a)
  =
  -\frac A{m-1}(x-a)^{1-m}+C.
\]
(3) 将 $x^2+px+q$ 进行配方得
\[
  x^2+px+q=\left(x+\frac p2\right)^2+\left(q-\frac{p^2}4\right),
\]
利用替换
\[
  z=x+\frac p2,\qquad x=z-\frac p2,\qquad \mathrm dx=\mathrm dz,
\]
又由于 $p^2-4q<0$, 有 $q-\dfrac{p^2}4>0$, 记 $a^2=q-\dfrac{p^2}4$, 得
\[
\begin{aligned}
  \int\frac{Ax+B}{x^2+px+q}\mathrm dx
  &=
  \int\frac{Ax+B}{\left(x+\dfrac p2\right)^2+\left(q-\dfrac{p^2}4\right)}\mathrm dx\\
  &=
  \int\frac{A(z-p/2)+B}{z^2+a^2}\mathrm dz\\
  &=
  \int\frac{Az}{z^2+a^2}\mathrm dz+
  \int\frac{B-Ap/2}{z^2+a^2}\mathrm dz\\
  &=
  \frac A2\ln(z^2+a^2)+\left(B-\frac{Ap}2\right)\frac1a\arctan\frac za+C\\
  &=
  \frac A2\ln(x^2+px+q)+
  \frac{2B-Ap}{\sqrt{4q-p^2}}\arctan\frac{2x+p}{\sqrt{4q-p^2}}+C.
\end{aligned}
\]

(4) 同样用变量替换 $z=x+\dfrac p2,\ x=z-\dfrac p2,\ \mathrm dx=\mathrm dz$, 且记 $a^2=q-\dfrac{p^2}4$, 得
\[
\begin{aligned}
  \int\frac{Ax+B}{(x^2+px+q)^n}\mathrm dx
  &=
  \int\frac{Az}{(z^2+a^2)^n}\mathrm dz+
  \left(B-\frac{Ap}2\right)\int\frac1{(z^2+a^2)^n}\mathrm dz.
\end{aligned}
\]
上式右边的第一项不定积分容易求出, 用换元积分法即可; 第二项积分要利用 \S 6.2 中例 6.2.26 的递推公式. 至此上述四种情况的不定积分全部得以解决。

现在我们讨论有理真分式的分解. 因为多项式 $Q(x)$ 在实数范围内能分解成一次因式和二次因式的乘积, 因此我们有
\[
  Q(x)=(x-a)^k\cdots(x-b)^t(x^2+px+q)^l\cdots(x^2+rx+s)^h,
\]
其中 $a,\ldots,b,p,q,\ldots,r,s$ 为常数; 且 $p^2-4q<0,\ldots,r^2-4s<0$; $k,\ldots,t,l,\ldots,h$ 为正整数. 可以证明真分式 $\dfrac{P(x)}{Q(x)}$ 必可分解成如下部分分式之和:
\[
\begin{aligned}
  \frac{P(x)}{Q(x)}
  &=
  \frac{A_1}{x-a}+\frac{A_2}{(x-a)^2}+\cdots+\frac{A_k}{(x-a)^k}+\cdots\\
  &\quad+
  \frac{B_1}{x-b}+\frac{B_2}{(x-b)^2}+\cdots+\frac{B_t}{(x-b)^t}\\
  &\quad+
  \frac{C_1x+D_1}{x^2+px+q}+\frac{C_2x+D_2}{(x^2+px+q)^2}
  +\cdots+\frac{C_lx+D_l}{(x^2+px+q)^l}+\cdots\\
  &\quad+
  \frac{E_1x+F_1}{x^2+rx+s}+\frac{E_2x+F_2}{(x^2+rx+s)^2}
  +\cdots+\frac{E_hx+F_h}{(x^2+rx+s)^h},
\end{aligned}
\]
其中 $A_1,A_2,\ldots,A_k; B_1,\ldots,B_t; C_1,\ldots,C_l; D_1,\ldots,D_l; E_1,\ldots,E_h; F_1,\ldots,F_h$ 都是常数. 因此有理函数的原函数都可由初等函数表示。

\noindent\textbf{例 6.3.1}\quad 求不定积分
\[
  I=\int\frac{x\mathrm dx}{(x^2+1)(x-1)}.
\]
\noindent\textbf{解}\quad 设
\[
  \frac{x}{(x^2+1)(x-1)}
  =
  \frac{Ax+B}{x^2+1}+\frac C{x-1}.
\]
将右端通分, 并比较等式两端分子同次幂的系数, 得
\[
  \begin{cases}
    A+C=0,\\
    -A+B=1,\\
    -B+C=0.
  \end{cases}
\]
解之得 $A=-\dfrac12,\ B=\dfrac12,\ C=\dfrac12$. 因此
\[
\begin{aligned}
  I
  &=
  -\frac12\int\frac{x-1}{x^2+1}\mathrm dx+\frac12\int\frac{\mathrm dx}{x-1}\\
  &=
  -\frac14\ln(x^2+1)+\frac12\arctan x+\frac12\ln|x-1|+C.
\end{aligned}
\]

\noindent\textbf{例 6.3.2}\quad 求不定积分
\[
  I=\int\frac{x^3+1}{x^4-3x^3+3x^2-x}\mathrm dx.
\]
\noindent\textbf{解}\quad 将分母作因式分解, 得
\[
  x^4-3x^3+3x^2-x=x(x-1)^3.
\]
于是, 我们可设
\[
  \frac{x^3+1}{x(x-1)^3}
  =
  \frac Ax+\frac B{(x-1)^3}+\frac C{(x-1)^2}+\frac D{x-1}.
\]
将两边乘 $x$, 然后令 $x=0$, 得 $A=-1$; 两边乘 $(x-1)^3$, 然后令 $x=1$, 得 $B=2$; 两边乘 $x-1$, 然后令 $x\to-\infty$, 得 $D=2$; 最后令 $x=-1$, 得 $C=1$. 这样, 我们便有
\[
\begin{aligned}
  I
  &=
  -\int\frac{\mathrm dx}x+2\int\frac{\mathrm dx}{(x-1)^3}
  +\int\frac{\mathrm dx}{(x-1)^2}+2\int\frac{\mathrm dx}{x-1}\\
  &=
  -\frac{x}{(x-1)^2}+\ln\frac{(x-1)^2}{|x|}+C.
\end{aligned}
\]

\subsection{三角函数有理式的不定积分}

二元有理函数是形如
\[
  R(u,v)=\frac{P(u,v)}{Q(u,v)}
\]
的两个变元的函数, 其中 $P(u,v)$ 和 $Q(u,v)$ 是二元多项式(即 $u^i,v^j,u^pv^q$ 的线性组合). 三角函数有理式是指形如 $R(\sin x,\cos x)$ 的函数, 其中 $R(u,v)$ 为二元有理函数, 它是由基本三角函数 $\sin x,\cos x,\tan x,\cot x$ 经有限次四则运算所得的函数. 但像如下形式的函数
\[
  \frac1{\sqrt{\sin^2x+5\cos^2x}},\qquad
  \sin x^2,\qquad
  \frac{\sin x}{x}
\]
则不属此列。

下面我们来说明三角函数有理式的不定积分
\[
  I=\int R(\sin x,\cos x)\mathrm dx
\]
一定可以表成初等函数. 令
\[
  t=\tan\frac x2\quad\text{或}\quad x=2\arctan t
\]
(通常称之为万能变换), 并且注意到
\[
\begin{aligned}
  \sin x
  &=
  2\sin\frac x2\cos\frac x2
  =
  \frac{2\tan\dfrac x2}{1+\tan^2\dfrac x2}
  =
  \frac{2t}{1+t^2},\\
  \cos x
  &=
  \cos^2\frac x2-\sin^2\frac x2
  =
  \frac{1-\tan^2\dfrac x2}{1+\tan^2\dfrac x2}
  =
  \frac{1-t^2}{1+t^2},\\
  \mathrm dx&=\frac{2\mathrm dt}{1+t^2},
\end{aligned}
\]
即有
\[
  I=
  \int R\left(\frac{2t}{1+t^2},\frac{1-t^2}{1+t^2}\right)
  \frac{2\mathrm dt}{1+t^2}.
\]
这样, 我们就将求三角函数有理式的不定积分变成了有理函数的不定积分. 再由上一小节所介绍的有理函数的积分方法就可解决求三角函数有理式的积分问题. 这也表明, 三角函数有理式的原函数一定可以表成初等函数。

\noindent\textbf{例 6.3.3}\quad 求不定积分
\[
  I=\int\frac{\mathrm dx}{5+4\sin x}.
\]
\noindent\textbf{解}\quad 令 $t=\tan\dfrac x2$, 则
\[
\begin{aligned}
  I
  &=
  \int\frac{\dfrac2{1+t^2}}{5+\dfrac{8t}{1+t^2}}\mathrm dt
  =
  \int\frac{2\mathrm dt}{5t^2+8t+5}\\
  &=
  \frac25\int\frac{\mathrm dt}{\left(t+\dfrac45\right)^2+\dfrac9{25}}
  =
  \frac23\arctan\frac{5t+4}3+C\\
  &=
  \frac23\arctan\frac{5\tan\dfrac x2+4}3+C.
\end{aligned}
\]

这里需指出的是, 具体解题时千万不要死搬硬套, 要具体问题作具体分析, 灵活运用. 事实上, 这种``万能变换''往往导致所要求的有理函数的积分比较复杂, 而对某些三角函数有理式的不定积分来说, 采用其他变换会更加方便. 例如, 当 $R(u,-v)=-R(u,v)$ 时, 我们可以证明, 这时必有 $R(u,v)=vR_1(u,v^2)$ 成立, 其中 $R_1$ 是一个二变元的有理函数. 于是若令 $t=\sin x$, 则有
\[
\begin{aligned}
  \int R(\sin x,\cos x)\mathrm dx
  &=
  \int\cos xR_1(\sin x,\cos^2x)\mathrm dx\\
  &=
  \int R_1(t,1-t^2)\mathrm dt.
\end{aligned}
\]
这表明, 若被积函数 $R(\sin x,\cos x)$ 关于 $\cos x$ 是奇函数, 则我们可用变换 $t=\sin x$ 将其转化为关于 $t$ 的有理函数的不定积分。

同理, 若 $R(-u,v)=-R(u,v)$, 则有 $R(u,v)=uR_1(u^2,v)$, 其中 $R_1$ 是一个二变元的有理函数. 于是若令 $t=\cos x$, 则有
\[
\begin{aligned}
  \int R(\sin x,\cos x)\mathrm dx
  &=
  \int\sin xR_1(\sin^2x,\cos x)\mathrm dx\\
  &=
  -\int R_1(1-t^2,t)\mathrm dt.
\end{aligned}
\]
这表明, 若被积函数 $R(\sin x,\cos x)$ 关于 $\sin x$ 是奇函数, 则我们可用变换 $t=\cos x$ 将其转化为关于 $t$ 的有理函数的不定积分。

再有若 $R(-u,-v)=R(u,v)$, 则可证此时有 $R(u,v)=R_1(u/v,v^2)$, 其中 $R_1$ 是一个二变元的有理函数. 于是, 若我们用变换 $t=\tan x$, 则有
\[
\begin{aligned}
  \int R(\sin x,\cos x)\mathrm dx
  &=
  \int R_1(\tan x,\cos^2x)\mathrm dx\\
  &=
  \int R_1\left(t,\frac1{1+t^2}\right)\frac{\mathrm dt}{1+t^2}.
\end{aligned}
\]
这表明, 若被积函数 $R(\sin x,\cos x)$ 关于 $\sin x$ 和 $\cos x$ 是偶函数, 则我们可用变换 $t=\tan x$ 将其转化为关于 $t$ 的有理函数的不定积分。

上述三种变换一般要比``万能变换''简单得多。

\noindent\textbf{例 6.3.4}\quad 求不定积分
\[
  I=\int\frac{\cos^3x}{1+\sin^2x}\mathrm dx.
\]
\noindent\textbf{解}\quad 由于 $R(\sin x,-\cos x)=-R(\sin x,\cos x)$, 故令 $t=\sin x$, 则
\[
\begin{aligned}
  I
  &=
  \int\frac{1-t^2}{1+t^2}\mathrm dt
  =
  \int\frac{2\mathrm dt}{1+t^2}-\int\mathrm dt\\
  &=
  2\arctan t-t+C
  =
  2\arctan\sin x-\sin x+C.
\end{aligned}
\]

\noindent\textbf{例 6.3.5}\quad 求不定积分
\[
  I=\int\frac{\sin x\cos x}{\sin^2x+\cos x}\mathrm dx.
\]
\noindent\textbf{解}\quad 由于 $R(-\sin x,\cos x)=-R(\sin x,\cos x)$, 故令 $t=\cos x$, 则
\[
\begin{aligned}
  I
  &=
  -\int\frac{\sin x\,\mathrm d\cos x}{\sin^2x+\cos x}
  =
  -\int\frac{t\mathrm dt}{1-t^2+t}\\
  &=
  \frac12\int\frac{\mathrm d(1+t-t^2)}{1+t-t^2}
  -\frac12\int\frac{\mathrm dt}{1+t-t^2}\\
  &=
  \frac12\ln|1+t-t^2|-\frac12\int\frac{\mathrm dt}{\dfrac54-\left(\dfrac12-t\right)^2}\\
  &=
  \frac12\ln|1+t-t^2|+
  \frac1{2\sqrt5}\ln\left|\frac{\sqrt5+1-2t}{\sqrt5-1+2t}\right|+C\\
  &=
  \frac12\left[
  \ln|1+\cos x-\cos^2x|+
  \frac1{\sqrt5}\ln\left|\frac{\sqrt5+1-2\cos x}{\sqrt5-1+2\cos x}\right|
  \right]+C.
\end{aligned}
\]

\noindent\textbf{例 6.3.6}\quad 求不定积分
\[
  I=\int\frac{\cos^2x}{(a^2\sin^2x+b^2\cos^2x)^2}\mathrm dx.
\]
\noindent\textbf{解}\quad 由于 $R(-\sin x,-\cos x)=R(\sin x,\cos x)$, 故令 $t=\tan x$, 则
\[
\begin{aligned}
  I
  &=
  \int\frac{\cos^4x}{(a^2\sin^2x+b^2\cos^2x)^2\cos^2x}\mathrm dx\\
  &=
  \int\frac{1}{(a^2\tan^2x+b^2)^2}\mathrm d\tan x\\
  &=
  \frac1{b^4}\int\frac{\mathrm dt}{\left(1+\dfrac{a^2}{b^2}t^2\right)^2}
  =
  \frac1{ab^3}\int\frac{\mathrm du}{(1+u^2)^2}\quad\left(u=\frac abt\right)\\
  &=
  \frac1{ab^3}\left\{\frac12\frac u{1+u^2}+\frac12\arctan u\right\}+C\\
  &=
  \frac{\tan x}{2b^2}\frac1{b^2+a^2\tan^2x}
  +\frac1{2ab^3}\arctan\left(\frac ab\tan x\right)+C.
\end{aligned}
\]

\noindent\textbf{例 6.3.7}\quad 求不定积分
\[
  I=\frac12\int\frac{1-r^2}{1-2r\cos x+r^2}\mathrm dx
  \qquad(0<r<1,\ |x|<\pi).
\]
\noindent\textbf{解}\quad 令 $\tan\dfrac x2=t$, 则有
\[
  \cos x=\frac{1-t^2}{1+t^2},\qquad
  \mathrm dx=\frac{2\mathrm dt}{1+t^2}.
\]
由此得
\[
\begin{aligned}
  I
  &=
  \int\frac{(1-r^2)\mathrm dt}{(1-r)^2+(1+r)^2t^2}\\
  &=
  \arctan\left(\frac{1+r}{1-r}t\right)+C\\
  &=
  \arctan\left(\frac{1+r}{1-r}\tan\frac x2\right)+C.
\end{aligned}
\]

\subsection{无理函数的不定积分}

这里我们考虑两类无理函数的不定积分. 第一类是如下的积分:
\[
  I=\int R\left(x,\sqrt[m]{\frac{ax+b}{cx+d}}\right)\mathrm dx,
\]
其中 $R(u,v)$ 是变元 $u,v$ 的二元函数, $m$ 为正整数, $a,b,c,d$ 为常数, 且 $ad-bc\ne0$. 令
\[
  \frac{ax+b}{cx+d}=t^m,
\]
则有
\[
  x=\frac{dt^m-b}{a-ct^m},\qquad
  \mathrm dx=\frac{m(ad-bc)t^{m-1}}{(a-ct^m)^2}\mathrm dt.
\]
于是, 我们有
\[
  I=
  \int R\left(\frac{dt^m-b}{a-ct^m},t\right)
  \frac{m(ad-bc)t^{m-1}}{(a-ct^m)^2}\mathrm dt.
\]
这样, 原不定积分就变成了有理函数的积分, 从而这类无理函数一定有初等函数作为其原函数。

\noindent\textbf{例 6.3.8}\quad 求不定积分
\[
  I=\int\frac{1-\sqrt{x+1}}{1+\sqrt[3]{x+1}}\mathrm dx.
\]
\noindent\textbf{解}\quad 令 $\sqrt[6]{x+1}=t$, 则有
\[
  I=6\int\frac{t^5-t^8}{1+t^2}\mathrm dt.
\]
设
\[
  \frac{t^5-t^8}{1+t^2}=P_6(t)+\frac{Bt+C}{1+t^2},
\]
其中 $P_6(t)$ 为 $t$ 的多项式, $B,C$ 是实数. 在上式两边乘以 $1+t^2$, 并且将 $t=\mathrm i=\sqrt{-1}$ 代入, 得
\[
  \mathrm i-1=B\mathrm i+C,
\]
于是有 $B=1,C=-1$, 而且
\[
\begin{aligned}
  P_6(t)
  &=
  \frac{t^5-t^8}{1+t^2}-\frac{t-1}{1+t^2}\\
  &=
  \frac{t^5-t^8-t+1}{1+t^2}
  =
  1-t-t^2+t^3+t^4-t^6.
\end{aligned}
\]
所以
\[
\begin{aligned}
  I
  &=
  6\int\left(1-t-t^2+t^3+t^4-t^6+\frac{t-1}{1+t^2}\right)\mathrm dt\\
  &=
  -\frac67t^7+\frac65t^5+\frac32t^4-2t^3-3t^2+6t
  +3\ln(1+t^2)-6\arctan t+C\\
  &=
  -\frac67(x+1)^{7/6}+\frac65(x+1)^{5/6}
  +\frac32(x+1)^{2/3}-2(x+1)^{1/2}-3(x+1)^{1/3}\\
  &\quad+
  6(x+1)^{1/6}+3\ln\left(1+(x+1)^{1/3}\right)
  -6\arctan(x+1)^{1/6}+C.
\end{aligned}
\]

我们要考虑的第二类无理函数的不定积分是如下的二项式微分式的积分:
\[
  I=\int x^m(a+bx^n)^p\mathrm dx,
\]
其中 $a,b$ 为常数, $m,n,p$ 为有理数. 令 $x^n=t$, 则 $\mathrm dx=\dfrac1n t^{1/n-1}\mathrm dt$, 从而有
\[
\begin{aligned}
  I
  &=
  \int t^{m/n}(a+bt)^p\frac1n t^{1/n-1}\mathrm dt\\
  &=
  \frac1n\int t^{(m+1)/n-1}(a+bt)^p\mathrm dt.
\end{aligned}
\]
为简明起见, 令 $q=\dfrac{m+1}n-1$, 则
\begin{equation}
  I=\int x^m(a+bx^n)^p\mathrm dx
  =
  \frac1n\int t^q(a+bt)^p\mathrm dt.
  \tag{6.3.1}
\end{equation}
当 $p$ 是整数而 $q=r/s(r,s$ 为整数) 时, 被积函数最多含有一个根式, 根式内的线性分式为 $t$, 因此只要作变换 $z=\sqrt[s]{t}=\sqrt[s]{x^n}$, 就可将原积分化为有理函数的积分。

当 $q$ 为整数而 $p=\mu/\lambda(\mu,\lambda$ 为整数) 时, 被积函数也最多含有一个根式, 根式内的线性分式为 $a+bt$, 因此只要作变换 $z=\sqrt[\lambda]{a+bt}=\sqrt[\lambda]{a+bx^n}$, 就可将原积分化为有理函数的积分。

最后, 因为积分 (6.3.1) 还可以改写成
\[
  I=\int x^m(a+bx^n)^p\mathrm dx
  =
  \frac1n\int t^{p+q}\left(\frac{a+bt}{t}\right)^p\mathrm dt,
\]
所以当 $p+q$ 是整数而 $p=\mu/\lambda(\mu,\lambda$ 为整数) 时, 被积函数最多含有一个根式, 根式内的线性分式为 $\dfrac{a+bt}{t}$, 因此只要作变换
\[
  z=\sqrt[\lambda]{\frac{a+bt}{t}}
  =
  \sqrt[\lambda]{\frac{a+bx^n}{x^n}},
\]
就可将原积分化为有理函数的积分。

总之, 如果在 $p,q$ 和 $p+q$ 之中至少有一个是整数, 则二项式微分式的积分就可以化成有理函数的积分, 从而也就可以表成初等函数。

19 世纪中叶, 切比雪夫曾经证明: 二项式微分式的积分, 除上述三种情形外, 都不能由初等函数表示. 比如, 积分
\[
  \int\frac{\mathrm dx}{\sqrt{1-x^4}}
\]
就不能用初等函数表示。

\noindent\textbf{例 6.3.9}\quad 求不定积分
\[
  I=\int\sqrt{x^2+\frac1{x^2}}\,\mathrm dx\qquad(x>0).
\]
\noindent\textbf{解}\quad 由于原积分可以写成
\[
  I=\int x^{-1}(1+x^4)^{1/2}\mathrm dx,
\]
因此
\[
  q=\frac{m+1}n-1=\frac{-1+1}4-1=-1
\]
为整数, 从而可用初等函数表示出. 具体计算如下:
\[
\begin{aligned}
  I
  &=
  \frac12\int\frac{\sqrt{1+x^4}}{x^2}\mathrm d(x^2)
  =
  \frac12\int\frac{\sqrt{1+t^2}}{t}\mathrm dt
  \qquad(x^2=t>0)\\
  &=
  \frac12\int\frac{1+t^2}{t\sqrt{1+t^2}}\mathrm dt\\
  &=
  \frac12\int\frac{t\mathrm dt}{\sqrt{1+t^2}}
  +\frac12\int\frac{\mathrm dt}{t^2\sqrt{1+1/t^2}}\\
  &=
  \frac12\sqrt{1+t^2}-\frac12\ln\left(\frac1t+\sqrt{1+\frac1{t^2}}\right)+C\\
  &=
  \frac12\sqrt{1+x^4}
  -\frac12\ln\left|\frac{1+\sqrt{1+x^4}}{x^2}\right|+C.
\end{aligned}
\]

\section*{习题六}

1. 证明: 若 $\displaystyle \int f(x)\mathrm dx=F(x)+C$, 则
\[
  \int f(ax+b)\mathrm dx=\frac1aF(ax+b)+C\qquad(a\ne0).
\]

2. 求下列不定积分:

(1) $\displaystyle \int\left(\sqrt{2x}+\sqrt[3]{3x}\right)\mathrm dx$;

(2) $\displaystyle \int\left(1+\frac1{x^2}\sqrt{x\sqrt x}\right)\mathrm dx$;

(3) $\displaystyle \int\left(3^{x+1}+3^{-x}+\frac13e^x\right)\mathrm dx$;

(4) $\displaystyle \int\frac{2-\sin^2x}{\cos^2x}\mathrm dx$;

(5) $\displaystyle \int\left(\sqrt{\frac{1+x}{1-x}}+\sqrt{\frac{1-x}{1+x}}\right)\mathrm dx$;

(6) $\displaystyle \int\frac{\cos2x}{\cos^2x\sin^2x}\mathrm dx$;

(7) $\displaystyle \int\sqrt{1-\sin2x}\,\mathrm dx$;

(8) $\displaystyle \int(x+|x|)\mathrm dx$;

(9) $\displaystyle \int\operatorname{sgn}(\sin x)\mathrm dx$;

(10) $\displaystyle \int[x]\mathrm dx$。

3. 用换元法求下列不定积分:

(1) $\displaystyle \int\frac{2\mathrm dx}{3-5x}$;

(2) $\displaystyle \int\frac{\mathrm dx}{\sqrt{2-x^2}}$;

(3) $\displaystyle \int x\sqrt[3]{1-3x}\,\mathrm dx$;

(4) $\displaystyle \int\frac{\mathrm dx}{1+\sin x}$;

(5) $\displaystyle \int x(1+x^2)^5\mathrm dx$;

(6) $\displaystyle \int\frac{x\mathrm dx}{\sqrt{1-x^2}}$;

(7) $\displaystyle \int\frac{x\mathrm dx}{4+x^4}$;

(8) $\displaystyle \int\frac{e^x}{1+e^x}\mathrm dx$;

(9) $\displaystyle \int\frac{\mathrm dx}{\sqrt x(1+x)}$;

(10) $\displaystyle \int\frac{\mathrm dx}{x\sqrt{x^2-1}}$;

(11) $\displaystyle \int\sin x\sin3x\,\mathrm dx$;

(12) $\displaystyle \int\sin^2x\cos^2x\,\mathrm dx$;

(13) $\displaystyle \int\frac{\sin2x}{\sqrt{1+\sin^2x}}\mathrm dx$;

(14) $\displaystyle \int\frac{\sin x\cos x}{\sqrt{a^2\sin^2x+b^2\cos^2x}}\mathrm dx\quad(a^2>b^2)$;

(15) $\displaystyle \int\frac{\sqrt{a^2-x^2}}x\mathrm dx$;

(16) $\displaystyle \int\frac{\mathrm dx}{x^2\sqrt{1+x^2}}$;

(17) $\displaystyle \int\frac{\mathrm dx}{\sqrt{(a^2+x^2)^3}}$;

(18) $\displaystyle \int\frac{x}{1+\sqrt x}\mathrm dx$;

(19) $\displaystyle \int\sqrt{\frac{2+x}{2-x}}\,\mathrm dx$;

(20) $\displaystyle \int\frac{\mathrm dx}{x\sqrt{4-x^2}}$。

4. 用分部积分法求下列不定积分:

(1) $\displaystyle \int\ln(1+x^2)\mathrm dx$;

(2) $\displaystyle \int xe^{-3x}\mathrm dx$;

(3) $\displaystyle \int x\cos nx\,\mathrm dx$;

(4) $\displaystyle \int\arcsin x\,\mathrm dx$;

(5) $\displaystyle \int\frac{x}{\sin^2x}\mathrm dx$;

(6) $\displaystyle \int\frac{\arctan x}{x^2(1+x^2)}\mathrm dx$;

(7) $\displaystyle \int\frac{e^{\arctan x}}{(1+x^2)^{3/2}}\mathrm dx$;

(8) $\displaystyle \int\frac{x^2}{(1+x^2)^2}\mathrm dx$;

(9) $\displaystyle \int\sin(\ln x)\mathrm dx$;

(10) $\displaystyle \int e^x\cos^2x\,\mathrm dx$。

5. 对下列不定积分建立递推公式:

(1) $\displaystyle I_n=\int x^m\ln^nx\,\mathrm dx$;

(2) $\displaystyle I_n=\int x^ne^{-x}\mathrm dx$;

(3) $\displaystyle I_n=\int\frac{x^n}{\sqrt{1-x^2}}\mathrm dx$;

(4) $\displaystyle I_n=\int\frac{\mathrm dx}{x^n\sqrt{1+x^2}}$。

6. 求下列有理函数的不定积分:

(1) $\displaystyle \int\frac{2x}{(x-1)(x+1)^2}\mathrm dx$;

(2) $\displaystyle \int\frac{2x^2+2x+13}{(x-2)(1+x^2)^2}\mathrm dx$;

(3) $\displaystyle \int\frac{x^2-x+1}{2x+1}\mathrm dx$;

(4) $\displaystyle \int\frac{3x+5}{(x^2+2x+2)^2}\mathrm dx$;

(5) $\displaystyle \int\frac{x+1}{x^3+2x^2-x-2}\mathrm dx$;

(6) $\displaystyle \int\frac{\mathrm dx}{x(x^3+2)}$;

(7) $\displaystyle \int\frac{\mathrm dx}{(x-2)^2(x+3)^3}$;

(8) $\displaystyle \int\frac{x\mathrm dx}{x^3-1}$。

7. 求下列三角函数有理式的不定积分:

(1) $\displaystyle \int\frac{\mathrm dx}{\sin^2x\cos x}$;

(2) $\displaystyle \int\frac{\mathrm dx}{\cos^2x\sin x}$;

(3) $\displaystyle \int\frac{\sin2x}{1+\cos^2x}\mathrm dx$;

(4) $\displaystyle \int\sin5x\sin3x\,\mathrm dx$;

(5) $\displaystyle \int\frac{\sin^2x}{1+\cos^2x}\mathrm dx$;

(6) $\displaystyle \int\frac{\sin^5x}{\cos^4x}\mathrm dx$;

(7) $\displaystyle \int\frac1{1+\cos x}\mathrm dx$;

(8) $\displaystyle \int\frac{1+\sin x}{1-\sin x}\mathrm dx$;

(9) $\displaystyle \int\frac{\sin x}{a\sin x+b\cos x}\mathrm dx$;

(10) $\displaystyle \int\frac{\sin x\mathrm dx}{\sin x-\cos x+2}$。

8. 求下列无理函数的不定积分:

(1) $\displaystyle \int\frac{\sqrt x}{\sqrt[4]{x^3}+1}\mathrm dx$;

(2) $\displaystyle \int\frac{\mathrm dx}{1+\sqrt x+\sqrt{x+1}}$;

(3) $\displaystyle \int\frac{\mathrm dx}{x^4\sqrt{1-x^2}}$;

(4) $\displaystyle \int\frac{\mathrm dx}{x\sqrt{x^2+x+1}}$;

(5) $\displaystyle \int\frac{\mathrm dx}{(x+1)\sqrt{x^2+4x+5}}$;

(6) $\displaystyle \int\frac{\mathrm dx}{x+\sqrt{x^2+x+1}}$;

(7) $\displaystyle \int\sqrt{\frac{1-x}{1+x}}\,\mathrm dx$;

(8) $\displaystyle \int\frac{\mathrm dx}{x\sqrt{4-x^2}}$;

(9) $\displaystyle \int\sqrt{\tan x}\,\mathrm dx$;

(10) $\displaystyle \int\sqrt[3]{\frac{2-x}{2+x}}\frac{\mathrm dx}{(2-x)^2}$。

\end{document}
